\subsection{LASSO and Elastic Net}
\label{sec:sparse_methods}
We used logistic-regression versions of LASSO and Elastic Net as sparse predictive baselines \cite{tibshirani1996lasso,zou2005elastic}. Both minimise a class-weighted logistic loss with a penalty on the gene coefficients $\boldsymbol{\beta}$ of the form
\begin{equation}
\Omega(\boldsymbol{\beta})=\alpha\lVert\boldsymbol{\beta}\rVert_1+
\frac{1-\alpha}{2}\lVert\boldsymbol{\beta}\rVert_2^2.
\label{eq:sparse_penalty}
\end{equation}
LASSO uses $\alpha=1$, whereas Elastic Net allows a mixture of $L_1$ and $L_2$ regularisation. Models were fitted using scikit-learn's SAGA solver with balanced class weights, at most 1,600 iterations, tolerance $10^{-3}$ and random seed 355. The inverse regularisation strength was searched over $C\in\{0.001,0.01,0.1,1\}$; Elastic Net additionally searched $\alpha\in\{0.5,0.9\}$.

For each outer LOCO split, we chose $(C,\alpha)$ using the mean AUROC across three inner cohort-held-out validation splits. Ties favoured smaller $C$, then larger $\alpha$. After refitting on all three development cohorts, genes with absolute standardised-model coefficients exceeding $10^{-8}$ formed the selected signature. The final model also predicted the untouched outer test cohort. Preprocessing was refitted within every inner training split as described in Section~\ref{sec:evaluation_design}; no fixed number of genes was imposed.

To evaluate gene agreement between independently analysed cohorts, we separately fitted LASSO and Elastic Net within each cohort. Here, hyperparameters were selected by stratified three-fold cross-validation within that cohort, followed by a fit on all its samples. We also recorded selection frequencies from 30 stratified half-sample refits per cohort and method as an exploratory sensitivity measure; these frequencies were not the formal Stability Selection procedure described separately. Within-cohort signatures were compared with each other, but their training-set predictions were not reported as external test results.
